\chapter{Rigorous finite-size seeds and an unconditional spectral gap}
\label{chap:seedkernel}

The spacetime-bootstrap criterion (Theorem~\ref{thm:rp-gap}) reduces an infinite-volume spectral
gap to the two finite-size purity inequalities~\eqref{eq:rp-seed-hyp} at a single seed scale
\((n_0, \beta_0)\):
\[
  \Ps(n_0, \beta_0) = \frac{Z_{2n_0}(\beta_0)}{Z_{n_0}(\beta_0)^2} \ge 1 - s_0, \qquad
  \Pt(2n_0, \beta_0) = \frac{Z_{2n_0}(2\beta_0)}{Z_{2n_0}(\beta_0)^2} \ge 1 - t_0.
\]
Certifying these inequalities unconditionally requires rigorous upper and lower bounds on the three
partition functions \(Z_{n_0}(\beta_0)\), \(Z_{2n_0}(\beta_0)\), and \(Z_{2n_0}(2\beta_0)\). At the
large inverse temperatures needed for the seed purities to approach \(1\), each partition function
\(Z_L(\beta) = \sum_{m=0}^{d^L-1} e^{-\beta E_L(m)}\) (Lemma~\ref{lem:ring-trace-exp}) is dominated
by the ground state, with a small correction from the first excited state and an exponentially
suppressed tail from all higher levels (Definition~\ref{def:low-spectrum} and
Lemma~\ref{lem:low-spectrum-z}). Rather than diagonalizing the full \(d^L \times d^L\) Hamiltonian,
it therefore suffices to:
\begin{enumerate}
  \item enclose the ground-state energy \(E_L(0)\) and first excited energy \(E_L(1)\) within narrow
  intervals using the Rayleigh--Ritz residuals of two approximate integer eigenvectors
  (Lemma~\ref{lem:residual}), and
  \item bound all remaining \(d^L - 2\) eigenvalues from below by a threshold \(\rho\) via a
  variational comparison with the diagonal part of the Hamiltonian (Lemmas~\ref{lem:counting}
  and~\ref{lem:ising-classical}).
\end{enumerate}
All matrix-vector products, norms, and residual bounds reduce to exact integer arithmetic
(Lemmas~\ref{lem:ising-bridge} and~\ref{lem:ising-low-spectra}). Carrying out this argument for the
mixed-field Ising chain~\eqref{eq:mixed-ising-ham} at \((g, h) = (1/4, 1/20)\) on rings of
\(n_0 = 6\) and \(2n_0 = 12\) sites proves the seed purity bounds (Theorem~\ref{thm:ising-seeds}),
discharging the finite-size seed hypothesis (Lemma~\ref{lem:ising-longitudinal-seeds}) of
Theorem~\ref{thm:ising-longitudinal-gap} at that point and yielding a fully unconditional spectral
gap theorem (Theorem~\ref{thm:ising-gap}).

\section{Perturbative and variational eigenvalue bounds}

We first record the two elementary spectral inequalities that turn finite-dimensional trial vectors
and quadratic-form comparisons into rigorous enclosures of the low-energy spectrum. Given an
approximate eigenvector \(x\) with Rayleigh quotient \(\theta\), the Euclidean norm of the residual
\(Ax - \theta x\) bounds the distance from \(\theta\) to the nearest eigenvalue of \(A\) via
Weinstein's inequality~\eqref{eq:rayleigh-residual}~\cite{Weinstein1934Modified}
(Lemma~\ref{lem:residual}). Complementarily, if the quadratic form
\(\operatorname{Re}\langle x, Ax\rangle\) is bounded below by \(\rho\|x\|^2\) on the subspace of
vectors vanishing on a finite set \(B\) of basis configurations, then at most \(|B|\) eigenvalues
of \(A\) can lie below \(\rho\) (Lemma~\ref{lem:counting}).

\begin{lemma}[Rayleigh--Ritz residual bound~\cite{Weinstein1934Modified}]
  \label{lem:residual}
  \lean{SeedKernel.Ising.exists_eigenvalue_near}
  \leanok
  For a Hermitian matrix \(A\), a vector \(x\), and \(\theta \in \R\), some eigenvalue \(\lambda_i\)
  of \(A\) satisfies
  \begin{equation}
    \label{eq:rayleigh-residual}
    (\lambda_i - \theta)^2 \|x\|^2 \le \|Ax - \theta x\|^2.
  \end{equation}
\end{lemma}
\begin{proof}
  \leanok
  Expanding in the unitary eigenbasis via \(y = U^\dagger x\) gives
  \(\|Ax - \theta x\|^2 = \sum_i (\lambda_i - \theta)^2 |y_i|^2 \ge
  \min_i (\lambda_i - \theta)^2 \|y\|^2\), where \(\|y\| = \|x\|\), which proves
  \eqref{eq:rayleigh-residual}.
\end{proof}

\begin{lemma}[Variational subspace eigenvalue bound]
  \label{lem:counting}
  \lean{SeedKernel.Ising.card_eigenvalues_lt_le}
  \leanok
  Let \(A\) be Hermitian, \(B\) a set of basis vectors, and \(\rho \in \R\). If
  \(\operatorname{Re}\langle x, Ax \rangle \ge \rho \|x\|^2\) for every vector \(x\) vanishing on
  \(B\), then at most \(|B|\) eigenvalues of \(A\) lie strictly below \(\rho\).
\end{lemma}
\begin{proof}
  \leanok
  If more than \(|B|\) eigenvalues were strictly below \(\rho\), their eigenvectors would span a
  subspace of dimension larger than \(|B|\), which by dimension counting must intersect the
  codimension-\(|B|\) subspace of vectors vanishing on \(B\) in a nonzero vector \(x\) satisfying
  \(\operatorname{Re}\langle x, Ax \rangle < \rho \|x\|^2\).
\end{proof}

\section{The mixed-field Ising Hamiltonian in the computational basis}

In the computational \(Z\)-basis, the mixed-field Ising Hamiltonian~\eqref{eq:mixed-ising-ham} acts
diagonally with the classical Ising energy \(E(\sigma)\) up to \(L\) single-spin flips of amplitude
\(g\) (Lemma~\ref{lem:ising-action}). By Cauchy--Schwarz (equivalently, Gershgorin's circle
theorem~\cite{Gershgorin1931Uber} applied to the off-diagonal spin-flip graph), the transverse
field shifts the quadratic form by at most \(gL\) in~\eqref{eq:ising-gershgorin}
(Lemma~\ref{lem:ising-classical}), so any configuration cutoff on the classical Ising energy
\(E(\sigma)\) immediately yields a variational lower bound on the quantum Hamiltonian via
Lemma~\ref{lem:counting}.

\begin{lemma}[Matrix action in the computational \(Z\)-basis]
  \label{lem:ising-action}
  \lean{SeedKernel.Ising.diagE, SeedKernel.Ising.flipAt, SeedKernel.Ising.hamiltonian_mulVec_apply}
  \leanok
  \uses{def:ising-longitudinal}
  For the mixed-field Ising ring~\eqref{eq:mixed-ising-ham} of
  Definition~\ref{def:ising-longitudinal}, with classical Ising energy
  \(E(\sigma) = -\sum_j z_j z_{j+1} - h \sum_j z_j\) (\(z_j = \pm 1\)),
  \((H_L x)(\sigma) = E(\sigma)\, x(\sigma) - g \sum_j x(\sigma^{(j)})\), where \(\sigma^{(j)}\)
  denotes the configuration obtained from \(\sigma\) by flipping spin \(j\).
\end{lemma}
\begin{proof}
  \leanok
  Evaluating the bond and on-site operators of Definition~\ref{def:ising-longitudinal} on
  computational basis configurations shows that \(Z \otimes Z\) and \(Z\) act diagonally with the
  classical Ising energy \(E(\sigma)\) while each \(X_j\) flips the \(j\)-th spin.
\end{proof}

\begin{lemma}[Gershgorin lower bound on the quadratic form]
  \label{lem:ising-classical}
  \lean{SeedKernel.Ising.re_quadratic_ge}
  \leanok
  \uses{lem:ising-action}
  For \(g \ge 0\),
  \begin{equation}
    \label{eq:ising-gershgorin}
    \operatorname{Re}\langle x, H_L x \rangle \ge \sum_\sigma \bigl(E(\sigma) - gL\bigr)\,
    |x(\sigma)|^2.
  \end{equation}
\end{lemma}
\begin{proof}
  \leanok
  Bounding each off-diagonal spin-flip product in Lemma~\ref{lem:ising-action} by
  \(\operatorname{Re}(\bar a b) \le (|a|^2 + |b|^2)/2\) and summing over the \(L\) single-spin flip
  neighbors of each configuration gives~\eqref{eq:ising-gershgorin}.
\end{proof}

\begin{lemma}[Binary encoding of spin configurations]
  \label{lem:ising-bridge}
  \lean{SeedKernel.Ising.encode, SeedKernel.Ising.applyH_encode, SeedKernel.Ising.dot_self_eq,
    SeedKernel.Ising.residualSq_eq}
  \leanok
  \uses{lem:ising-action}
  When spin configurations are encoded as the binary digits of an integer in
  \(\{0, \dots, 2^L - 1\}\) so that flipping spin \(j\) toggles the \(j\)-th bit (bitwise
  exclusive-or with \(2^j\)), the integer action of \(20 H_L\) and the bit-indexed sums equal the
  matrix action of \(H_L\) in Lemma~\ref{lem:ising-action}, the squared norm, and the squared
  residual of the decoded vector.
\end{lemma}
\begin{proof}
  \leanok
  Rewrite the sums over \(\{0, \dots, 2^L - 1\}\) as sums over spin configurations via the binary
  encoding bijection.
\end{proof}

\section{Certified seed purities and the unconditional spectral gap}

Combining the two Rayleigh--Ritz residual bounds (Lemma~\ref{lem:residual}) with the Gershgorin
comparison (Lemmas~\ref{lem:counting} and~\ref{lem:ising-classical}) yields the low-energy spectral
enclosures of Lemma~\ref{lem:ising-low-spectra} and the resulting two-sided
bounds~\eqref{eq:low-spectrum-z} on the partition functions \(Z_6(97/4)\), \(Z_{12}(97/4)\), and
\(Z_{12}(97/2)\). These bounds certify the seed purities~\eqref{eq:ising-certified-seeds} at
\((g, h) = (1/4, 1/20)\) (Theorem~\ref{thm:ising-seeds}) and establish the unconditional spectral
gap~\eqref{eq:ising-uncond-gap} for all even ring lengths \(L \ge 6\) (Theorem~\ref{thm:ising-gap}).

\begin{definition}[Low-energy spectral enclosure]
  \label{def:low-spectrum}
  \lean{SeedKernel.Ising.LowSpectrum}
  \leanok
  An eigenvalue family \(e\) admits a \emph{low-energy spectral enclosure} \((a_0, b_0, a_1, \rho)\)
  if its lowest eigenvalue lies in \([a_0, b_0]\), its second eigenvalue is at least \(a_1\), and
  all remaining eigenvalues are at least \(\rho\).
\end{definition}

\begin{lemma}[Partition-function bounds from a spectral enclosure]
  \label{lem:low-spectrum-z}
  \lean{SeedKernel.Ising.sum_exp_ge, SeedKernel.Ising.sum_exp_le}
  \leanok
  \uses{def:low-spectrum}
  For \(\beta \ge 0\) and a spectrum of \(N\) eigenvalues admitting the enclosure
  \((a_0, b_0, a_1, \rho)\) of Definition~\ref{def:low-spectrum},
  \begin{equation}
    \label{eq:low-spectrum-z}
    e^{-\beta b_0} \le \sum_{i=0}^{N-1} e^{-\beta e_i} \le e^{-\beta a_0} + e^{-\beta a_1} +
    N e^{-\beta \rho}.
  \end{equation}
\end{lemma}
\begin{proof}
  \leanok
  Dropping all Boltzmann weights except that of the ground state in \([a_0, b_0]\) gives the lower
  bound in~\eqref{eq:low-spectrum-z}, and bounding each Boltzmann weight by the corresponding
  spectral lower bound of Definition~\ref{def:low-spectrum} gives the upper bound.
\end{proof}

\begin{lemma}[Spectral enclosures of the \(6\)- and \(12\)-site Ising rings]
  \label{lem:ising-low-spectra}
  \lean{SeedKernel.Ising.lowSpectrum_of_checks, SeedKernel.Ising.lowSpectrum6,
    SeedKernel.Ising.lowSpectrum12, SeedKernel.Ising.allChunks12_eq_range}
  \leanok
  \uses{lem:residual, lem:counting, lem:ising-classical, lem:ising-bridge, def:low-spectrum}
  At \(g = 1/4\) and \(h = 1/20\), the rings of \(6\) and \(12\) sites admit low-energy spectral
  enclosures (Definition~\ref{def:low-spectrum}) with eigenvalue radius \(10^{-10}\) around the two
  Ritz values of the integer trial vectors and excited-state threshold \(\rho = c/20 - L/4\), where
  exactly two classical configurations satisfy \(20 E(\sigma) < c\).
\end{lemma}
\begin{proof}
  \leanok
  Exact integer arithmetic via Lemma~\ref{lem:ising-bridge} verifies the norms and residuals of two
  integer trial vectors on each ring (evaluated for \(12\) sites in \(64\) blocks whose union
  covers \(\{0, \dots, 4095\}\)) together with the count of classical configurations below the
  energy cutoff. Lemma~\ref{lem:residual} converts the residual bounds~\eqref{eq:rayleigh-residual}
  into eigenvalue enclosures, and Lemmas~\ref{lem:counting} and~\ref{lem:ising-classical} convert
  the classical configuration count in~\eqref{eq:ising-gershgorin} into the lower bound \(\rho\) on
  all remaining eigenvalues.
\end{proof}

\begin{theorem}[Certified seed purities of the mixed-field Ising ring]
  \label{thm:ising-seeds}
  \lean{SeedKernel.Ising.seed_spatial, SeedKernel.Ising.seed_thermal,
    SeedKernel.Ising.first_point_seed}
  \leanok
  \uses{lem:ising-low-spectra, lem:low-spectrum-z, def:rp-purities, def:ising-longitudinal}
  For \(H_L = -\sum_j Z_j Z_{j+1} - \tfrac14 \sum_j X_j - \tfrac1{20} \sum_j Z_j\),
  \begin{equation}
    \label{eq:ising-certified-seeds}
    \Ps(6, 97/4) \ge 1 - 9064 \cdot 10^{-9}, \qquad \Pt(12, 97/4) \ge 1 - 3141 \cdot 10^{-8}.
  \end{equation}
  These are the seed purity inequalities~\eqref{eq:rp-seed-hyp} of the first grid point of
  Definition~\ref{def:ising-longitudinal} (Lemma~\ref{lem:ising-longitudinal-seeds}).
\end{theorem}
\begin{proof}
  \leanok
  Applying~\eqref{eq:low-spectrum-z} from Lemma~\ref{lem:low-spectrum-z} to the spectral enclosures
  of Lemma~\ref{lem:ising-low-spectra} gives \(Z \ge e^{-\beta(\theta_0 + r)}\) and
  \(Z \le e^{-\beta(\theta_0 - r)}(1 + T)\) with an exponentially small excited-state sum \(T\);
  the purity ratios \(Z_{12}(\beta)/Z_6(\beta)^2\) and \(Z_{12}(2\beta)/Z_{12}(\beta)^2\) in
  Definition~\ref{def:rp-purities} are then bounded from below by \((1 - x)/(1 + t)^2\) using the
  rational bound \(e^{-14} \le (27182818283/10^{10})^{-14}\), proving~\eqref{eq:ising-certified-seeds}.
\end{proof}

\begin{theorem}[Unconditional spectral gap of the mixed-field Ising ring]
  \label{thm:ising-gap}
  \lean{SeedKernel.Ising.firstPoint_mem, SeedKernel.Ising.gap_first_point,
    SeedKernel.Ising.ising_gap}
  \leanok
  \uses{thm:ising-seeds, lem:ising-longitudinal-valid, thm:rp-gap}
  For every even \(L \ge 6\), the mixed-field Ising ring
  \(-\sum_j Z_j Z_{j+1} - \tfrac14 \sum_j X_j - \tfrac1{20} \sum_j Z_j\) has a unique ground state
  and
  \begin{equation}
    \label{eq:ising-uncond-gap}
    \Delta_L = E_1(L) - E_0(L) > \frac{\log 13480}{97/4} \approx 0.3921 .
  \end{equation}
\end{theorem}
\begin{proof}
  \leanok
  Apply Theorem~\ref{thm:rp-gap} at the first grid point of Definition~\ref{def:ising-longitudinal},
  using the certificate admissibility from Lemma~\ref{lem:ising-longitudinal-valid} and the seed
  purity bounds~\eqref{eq:ising-certified-seeds} from Theorem~\ref{thm:ising-seeds}. This discharges
  the seed hypothesis of Lemma~\ref{lem:ising-longitudinal-seeds} at \((g, h) = (1/4, 1/20)\),
  establishing the spectral gap~\eqref{eq:ising-uncond-gap} of Theorem~\ref{thm:ising-longitudinal-gap}
  unconditionally.
\end{proof}

\section{Exact integer arithmetic for spin-1 trial states}

\begin{remark}[Rayleigh--Ritz residuals for filtered spin-1 trial states]
  \label{rem:spin1-trial}
  The Rayleigh--Ritz residual strategy of Lemma~\ref{lem:residual} extends to integer trial states
  of the spin-1 Heisenberg chain~\eqref{eq:heisenberg-ham}. Applying a shifted polynomial filter
  \(p(H_n)\) to a N\'eel product state in exact \(3^n\)-dimensional integer arithmetic yields
  certified ground-state energy upper bounds and residual enclosures~\eqref{eq:rayleigh-residual}
  on rings of \(n = 6, 8, 10\) sites. Unlike the weakly perturbed Ising chain of
  Theorem~\ref{thm:ising-gap}, however, the isotropic spin-1 Heisenberg chain is strongly coupled
  and requires a larger seed ring (\(n_0 = 36\) in Lemma~\ref{lem:heisenberg-seed36}), where a
  simple diagonal Gershgorin cutoff such as~\eqref{eq:ising-gershgorin} does not suffice to bound
  the remaining excited spectrum; see~\cite{OpenAIHaldane} for the twisted-ring trace recursion
  and polynomial-filter method used to certify the spin-1 seeds~\eqref{eq:heisenberg-seed36}.
\end{remark}

